\begin{frame}{Proofs by Cases}
\begin{center}
\vspace{1.5cm}
{\Huge\textbf{Proofs by Cases}}
\end{center}
\end{frame}

% --------------------------------------------------
% Question 1
% --------------------------------------------------

\begin{frame}{Question 1}
A software program accepts any integer input $n$ and computes:

$$
f(n)=|n|
$$

Show that:

$$
|n|\geq 0
$$

for every integer $n$ by considering the possible cases for $n$.
\end{frame}

\begin{frame}{Question 1 -- Answer}
The two possible cases are:

\begin{itemize}
\item $n\geq 0$
\item $n<0$
\end{itemize}

We examine the absolute value separately in each case.
\end{frame}

\begin{frame}{Question 1 -- Answer}
\textbf{Case 1: $n\geq0$}

By definition:

$$
|n|=n
$$

Since $n\geq0$:

$$
|n|\geq0
$$

\end{frame}

\begin{frame}{Question 1 -- Answer}
\textbf{Case 2: $n<0$}

By definition:

$$
|n|=-n
$$

Since $n<0$:

$$
-n>0
$$

Therefore:

$$
|n|\geq0
$$

Hence, in both cases:

$$
\boxed{|n|\geq0}
$$

\end{frame}

% --------------------------------------------------
% Question 2
% --------------------------------------------------

\begin{frame}{Question 2}
A hashing algorithm produces a hash value using:

$$
h(n)=n^2
$$

For an integer input $n$, determine whether the hash value has the same
parity as $n$.

Prove your conclusion by considering the cases where $n$ is even and
where $n$ is odd.
\end{frame}

\begin{frame}{Question 2 -- Answer}
There are two possible cases.

\begin{itemize}
\item $n$ is even
\item $n$ is odd
\end{itemize}

We examine the parity of $n^2$ in each case.
\end{frame}

\begin{frame}{Question 2 -- Answer}
\textbf{Case 1: $n$ is even}

Let:

$$
n=2k
$$

Then:

$$
n^2=(2k)^2
$$

$$
n^2=4k^2=2(2k^2)
$$

Therefore, $n^2$ is even.
\end{frame}

\begin{frame}{Question 2 -- Answer}
\textbf{Case 2: $n$ is odd}

Let:

$$
n=2k+1
$$

Then:

$$
n^2=(2k+1)^2
$$

$$
n^2=4k^2+4k+1
$$

$$
n^2=2(2k^2+2k)+1
$$

Therefore, $n^2$ is odd.

Hence, $n^2$ has the same parity as $n$.
\end{frame}

\begin{frame}{Question 2 -- Answer}
Therefore:

$$
\boxed{n^2\text{ has the same parity as }n}
$$

for every integer $n$.
\end{frame}

% --------------------------------------------------
% Question 3
% --------------------------------------------------

\begin{frame}{Question 3}
A scheduling system assigns an integer priority value $n$ to a task.
The system calculates:

$$
g(n)=n(n+1)
$$

Show that the calculated value is always even by considering the two
possible parity cases for $n$.
\end{frame}

\begin{frame}{Question 3 -- Answer}
There are two cases:

\begin{itemize}
\item $n$ is even
\item $n$ is odd
\end{itemize}

We prove that $n(n+1)$ is even in both cases.
\end{frame}

\begin{frame}{Question 3 -- Answer}
\textbf{Case 1: $n$ is even}

Let:

$$
n=2k
$$

Then:

$$
n(n+1)=2k(2k+1)
$$

Therefore:

$$
n(n+1)=2[k(2k+1)]
$$

Hence, the result is even.
\end{frame}

\begin{frame}{Question 3 -- Answer}
\textbf{Case 2: $n$ is odd}

Let:

$$
n=2k+1
$$

Then:

$$
n(n+1)=(2k+1)(2k+2)
$$

Therefore:

$$
n(n+1)=2[(2k+1)(k+1)]
$$

Hence, the result is also even.
\end{frame}

\begin{frame}{Question 3 -- Answer}
Since both possible cases produce an even result:

$$
\boxed{n(n+1)\text{ is always even}}
$$

\end{frame}

% --------------------------------------------------
% Question 4
% --------------------------------------------------

\begin{frame}{Question 4}
A computer program compares two integers $a$ and $b$ and returns the
larger value:

$$
\max(a,b)
$$

Prove that:

$$
\max(a,b)\geq a
$$

for all integers $a$ and $b$ by considering the possible relationships
between $a$ and $b$.
\end{frame}

\begin{frame}{Question 4 -- Answer}
There are two cases:

\begin{itemize}
\item $a\geq b$
\item $a<b$
\end{itemize}

We consider the value returned by the program in each case.
\end{frame}

\begin{frame}{Question 4 -- Answer}
\textbf{Case 1: $a\geq b$}

In this case:

$$
\max(a,b)=a
$$

Therefore:

$$
\max(a,b)\geq a
$$

because equality holds.
\end{frame}

\begin{frame}{Question 4 -- Answer}
\textbf{Case 2: $a<b$}

In this case:

$$
\max(a,b)=b
$$

Since:

$$
a<b
$$

we have:

$$
b>a
$$

Therefore:

$$
\max(a,b)>a
$$

\end{frame}

\begin{frame}{Question 4 -- Answer}
Both possible cases satisfy the required inequality.

Hence:

$$
\boxed{\max(a,b)\geq a}
$$

for all integers $a$ and $b$.
\end{frame}

% --------------------------------------------------
% Question 5
% --------------------------------------------------

\begin{frame}{Question 5}
A data-processing algorithm assigns a cost to an integer input $n$ using:

$$
C(n)=
\begin{cases}
n+2, & n\geq0\\
-n+2, & n<0
\end{cases}
$$

Show that:

$$
C(n)\geq2
$$

for every integer $n$ by proving the result separately for the two
possible cases.
\end{frame}

\begin{frame}{Question 5 -- Answer}
The definition of $C(n)$ naturally gives two cases:

\begin{itemize}
\item $n\geq0$
\item $n<0$
\end{itemize}

We prove the required inequality in each case.
\end{frame}

\begin{frame}{Question 5 -- Answer}
\textbf{Case 1: $n\geq0$}

The function is:

$$
C(n)=n+2
$$

Since:

$$
n\geq0
$$

we obtain:

$$
n+2\geq2
$$

Therefore:

$$
C(n)\geq2
$$

\end{frame}

\begin{frame}{Question 5 -- Answer}
\textbf{Case 2: $n<0$}

The function is:
$$
C(n)=-n+2
$$

Since:
$$
n<0
$$

we have:
$$
-n>0
$$

Therefore:
$$
-n+2>2
$$

and hence:
$$
C(n)\geq2
$$

\end{frame}

\begin{frame}{Question 5 -- Answer}
Both possible cases satisfy the required condition.

Therefore:

$$
\boxed{C(n)\geq2}
$$

for every integer $n$.

The result follows by examining all cases defined by the function.
\end{frame}
